\chapter{Escopo científico e limitações}

\eyebrow{o que os números podem e não podem afirmar}

\vspace{4pt}
O Apex foi construído para ser honesto quanto à fronteira entre o que suas
medições estabelecem e o que permanece em aberto. Este capítulo declara essa
fronteira de forma explícita.

\section{O que é validado}

\begin{itemize}[leftmargin=1.3em,itemsep=3pt]
  \item A \textbf{segmentação de gordura} é validada (IoU $\sim$0,92) e confiável,
        desde que o fundo tenha sido removido para que o modelo veja as condições
        em que foi treinado.
  \item O \textbf{pareamento verificável} é garantido por construção: toda imagem
        carrega seu \tele{carcass\_id} desde o momento da captura ou importação.
  \item A \textbf{convexidade de conformação} é uma medição objetiva e
        reproduzível do perfil, invariante à pose e à iluminação; sua
        correspondência com a anatomia é verificável a olho nu.
\end{itemize}

\section{O que não é validado}

Os dois \emph{graus} --- acabamento e conformação --- são estimativas, não
predições estabelecidas. Isso se fundamenta no experimento oráculo do projeto,
que testou se a conformação é recuperável a partir de uma única imagem 2D:

\begin{caution}[title=O resultado do oráculo]
Calculando descritores de forma sobre a \emph{anotação manual perfeita} --- a
melhor máscara que qualquer modelo poderia produzir --- \textbf{zero de 384
testes estatísticos sobreviveram à correção de múltiplas comparações} com n=22.
Mesmo com segmentação perfeita, a conformação não é recuperável de forma
confiável a partir de uma única imagem 2D neste tamanho de amostra. O gargalo é
a modalidade de aquisição, não o modelo.
\end{caution}

Portanto, o software apresenta esses graus como \pillalert{estimativas}, nunca
como fatos, e os marca em âmbar onde quer que apareçam.

\section{Por que isso importa --- e como se resolve}

A validação que falta aos graus requer exatamente uma coisa: graus corretamente
pareados às imagens, em uma amostra maior. É precisamente isso que o Apex foi
construído para coletar. À medida que o conjunto de dados pareado cresce, a aba
\textbf{Correlação física} (Estação~7) pode testar se as medições de superfície
predizem a referência física e os graus dos avaliadores --- fechando o ciclo que
a pesquisa anterior não conseguiu.

\section{Ressalvas de modalidade}

\begin{itemize}[leftmargin=1.3em,itemsep=3pt]
  \item A \textbf{cobertura de gordura (acabamento)} é uma propriedade de
        profundidade; uma única imagem de superfície a captura apenas
        parcialmente. A captura de profundidade (Kinect) é suportada no esquema
        para trabalhos futuros.
  \item A \textbf{captura padronizada} importa: um fundo fixo, iluminação uniforme
        e uma vista consistente melhoram toda medição subsequente. O software
        impõe uma vista obrigatória e oferece a captura de fundo, mas a montagem
        física é responsabilidade do operador.
  \item \textbf{Tamanho da amostra.} Os modelos atuais vêm de n=22. Suas
        estimativas devem ser lidas como exploratórias até que a amostra alvo seja
        coletada e o ciclo de validação seja executado.
\end{itemize}

\begin{note}[title=O fio condutor]
O estudo anterior não falhou em seus métodos --- falhou na integridade dos dados.
O Apex existe para corrigir isso primeiro. Cada limitação honesta acima é uma
razão pela qual o software é construído do jeito que é: coletar um conjunto de
dados pareado e confiável, e deixar a validação seguir.
\end{note}
